
\documentclass{article}
\usepackage{colortbl}
\usepackage{makecell}
\usepackage{multirow}
\usepackage{supertabular}

\begin{document}

\newcounter{utterance}

\centering \large Interaction Transcript for game `eqbench', experiment `eq\_bench', episode 23 with reflexion{-}llama{-}3.1{-}8b{-}instruct.
\vspace{24pt}

{ \footnotesize  \setcounter{utterance}{1}
\setlength{\tabcolsep}{0pt}
\begin{supertabular}{c@{$\;$}|p{.15\linewidth}@{}p{.15\linewidth}p{.15\linewidth}p{.15\linewidth}p{.15\linewidth}p{.15\linewidth}}
    \# & \multicolumn{2}{c}{Player} && \multicolumn{2}{c}{Game Master} \\
    \hline

    \theutterance \stepcounter{utterance}  
    & & & \multicolumn{4}{p{0.6\linewidth}}{
        \cellcolor[rgb]{0.9,0.9,0.9}{
            \makecell[{{p{\linewidth}}}]{
                \texttt{\tiny{[P1$\langle$GM]}}
                \texttt{You have attempted this same task before and did not succeed. These are the plans you wrote after those attempts.} \\
\texttt{Trial \#1: New plan: I got stuck in a simplistic analysis of the emotions, focusing on the surface{-}level reactions of Nox and Eva. However, I failed to consider the underlying dynamics of their relationship and the context of the dialogue. Next time, I will delve deeper into the emotional nuances of the conversation, taking into account the power dynamics, past conflicts, and the character's motivations. I will also consider the emotional resonance of the dialogue, such as the use of sarcasm, tone, and body language, to better understand Nox's emotional state. Specifically, I will analyze the following:} \\
\\ 
\texttt{* Nox's statement "I'm not playing anything, Eva. I'm just trying to move forward" as a potential attempt to deflect responsibility or avoid confrontation.} \\
\texttt{* Eva's accusation "You're still playing the victim" as a possible trigger for Nox's emotional response.} \\
\texttt{* The context of Nox's "studio" and his "producing for kids" as a potential source of guilt or shame.} \\
\\ 
\texttt{By considering these factors, I will aim to provide a more nuanced and accurate assessment of Nox's emotional state.} \\
\\ 
\texttt{Trial \#2: New plan: I got stuck in a simplistic analysis of the emotions, focusing on the surface{-}level reactions of Nox and Eva. However, I failed to consider the underlying dynamics of their relationship and the context of the dialogue. Next time, I will delve deeper into the emotional nuances of the conversation, taking into account the power dynamics, past conflicts, and the character's motivations. I will also consider the emotional resonance of the dialogue, such as the use of sarcasm, tone, and body language, to better understand Nox's emotional state. Specifically, I will analyze the following:} \\
\\ 
\texttt{* Nox's statement "I'm not playing anything, Eva. I'm just trying to move forward" as a potential attempt to deflect responsibility or avoid confrontation.} \\
\texttt{* Eva's accusation "You're still playing the victim" as a possible trigger for Nox's emotional response.} \\
\texttt{* The context of Nox's "studio" and his "producing for kids" as a potential source of guilt or shame.} \\
\texttt{* The fact that Eva is accusing Nox of hiding in his studio, which could imply that Nox is avoiding confrontation or responsibility, and this could be a source of guilt or shame for him.} \\
\texttt{* The tone of Eva's accusations, which could be perceived as condescending or mocking, and this could further trigger Nox's emotional response.} \\
\\ 
\texttt{By considering these factors, I will aim to provide a more nuanced and accurate assessment of Nox's emotional state.} \\
\\ 
\texttt{Trial \#3: New plan: I got stuck in a simplistic analysis of the emotions, focusing on the surface{-}level reactions of Nox and Eva. However, I failed to consider the underlying dynamics of their relationship and the context of the dialogue. Next time, I will delve deeper into the emotional nuances of the conversation, taking into account the power dynamics, past conflicts, and the character's motivations. I will also consider the emotional resonance of the dialogue, such as the use of sarcasm, tone, and body language, to better understand Nox's emotional state. Specifically, I will analyze the following:} \\
\\ 
\texttt{* Nox's statement "I'm not playing anything, Eva. I'm just trying to move forward" as a potential attempt to deflect responsibility or avoid confrontation.} \\
\texttt{* Eva's accusation "You're still playing the victim" as a possible trigger for Nox's emotional response.} \\
\texttt{* The context of Nox's "studio" and his "producing for kids" as a potential source of guilt or shame.} \\
\texttt{* The fact that Eva is accusing Nox of hiding in his studio, which could imply that Nox is avoiding confrontation or responsibility, and this could be a source of guilt or shame for him.} \\
\texttt{* The tone of Eva's accusations, which could be perceived as condescending or mocking, and this could further trigger Nox's emotional response.} \\
\texttt{* I will also consider the fact that Nox's statement "Not after everything" suggests a past conflict or trauma, which could be a significant factor in his emotional response.} \\
\texttt{* I will analyze the emotional tone of Nox's statement "I'm just trying to move forward" and consider whether it is a genuine attempt to move on or a defensive mechanism to avoid confronting his past.} \\
\\ 
\texttt{By considering these factors, I will aim to provide a more nuanced and accurate assessment of Nox's emotional state.} \\
\\ 
\\ 
\\ 
\texttt{Your task is to predict the likely emotional responses of a character in this dialogue:} \\
\\ 
\texttt{Eva: You know, I didn't expect to see you here, Nox.} \\
\texttt{Nox: And I didn't expect you to show up, Eva. Not after everything.} \\
\texttt{Eva: Oh, so you're still playing the victim, huh?} \\
\texttt{Nox: I'm not playing anything, Eva. I'm just trying to move forward.} \\
\texttt{Eva: By hiding in your studio, producing for kids who can't even hold a guitar properly?} \\
\texttt{{[}End dialogue{]}} \\
\\ 
\texttt{At the end of this dialogue, Nox would feel...} \\
\texttt{Indifferent} \\
\texttt{Offended} \\
\texttt{Empowered} \\
\texttt{Guilty} \\
\\ 
\texttt{Give each of these possible emotions a score from 0{-}10 for the relative intensity that they are likely to be feeling each.} \\
\\ 
\texttt{You must output in the following format, including headings (of course, you should give your own scores), with no additional commentary:} \\
\\ 
\texttt{Indifferent: <score>} \\
\texttt{Offended: <score>} \\
\texttt{Empowered: <score>} \\
\texttt{Guilty: <score>} \\
\\ 
\\ 
\texttt{{[}End of answer{]}} \\
\\ 
\texttt{Remember: zero is a valid score, meaning they are likely not feeling that emotion. You must score at least one emotion > 0.} \\
\\ 
\texttt{Your answer:} \\
            }
        }
    }
     \\ \\

    \theutterance \stepcounter{utterance}  
    & \multicolumn{4}{p{0.6\linewidth}}{
        \cellcolor[rgb]{0.561,0.737,0.561}{
            \makecell[{{p{\linewidth}}}]{
                \texttt{\tiny{[P1$\rangle$GM]}}
                \texttt{Indifferent: 0} \\
\texttt{Offended: 8} \\
\texttt{Empowered: 0} \\
\texttt{Guilty: 9} \\
            }
        }
    }
    & & \\ \\

    \theutterance \stepcounter{utterance}  
    & & & \multicolumn{2}{p{0.3\linewidth}}{
        \cellcolor[rgb]{0.9,0.9,0.9}{
            \makecell[{{p{\linewidth}}}]{
                \texttt{\tiny{[GM$|$GM]}}
                \texttt{Indifferent: 0} \\
\texttt{Offended: 8} \\
\texttt{Empowered: 0} \\
\texttt{Guilty: 9} \\
            }
        }
    }
    & & \\ \\

    \theutterance \stepcounter{utterance}  
    & & & \multicolumn{2}{p{0.3\linewidth}}{
        \cellcolor[rgb]{0.9,0.9,0.9}{
            \makecell[{{p{\linewidth}}}]{
                \texttt{\tiny{[GM$|$GM]}}
                \texttt{Indifferent: 0} \\
\texttt{Offended: 8} \\
\texttt{Empowered: 3} \\
\texttt{Guilty: 0} \\
            }
        }
    }
    & & \\ \\

    \theutterance \stepcounter{utterance}  
    & & & \multicolumn{2}{p{0.3\linewidth}}{
        \cellcolor[rgb]{0.9,0.9,0.9}{
            \makecell[{{p{\linewidth}}}]{
                \texttt{\tiny{[GM$|$GM]}}
                \texttt{game\_result = LOSE} \\
            }
        }
    }
    & & \\ \\

    \theutterance \stepcounter{utterance}  
    & & \multicolumn{4}{p{0.6\linewidth}}{
        \cellcolor[rgb]{0.95,0.95,0.95}{
            \makecell[{{p{\linewidth}}}]{
                \texttt{\tiny{[P1$\leftrightarrow$P]}}
                \texttt{Post{-}episode Reflexion (not sent to GM):} \\
\texttt{New plan: I got stuck in a simplistic analysis of the emotions, focusing on the surface{-}level reactions of Nox and Eva. However, I failed to consider the underlying dynamics of their relationship and the context of the dialogue. Next time, I will delve deeper into the emotional nuances of the conversation, taking into account the power dynamics, past conflicts, and the character's motivations. I will also consider the emotional resonance of the dialogue, such as the use of sarcasm, tone, and body language, to better understand Nox's emotional state. Specifically, I will analyze the following:} \\
\\ 
\texttt{* Nox's statement "I'm not playing anything, Eva. I'm just trying to move forward" as a potential attempt to deflect responsibility or avoid confrontation.} \\
\texttt{* Eva's accusation "You're still playing the victim" as a possible trigger for Nox's emotional response.} \\
\texttt{* The context of Nox's "studio" and his "producing for kids" as a potential source of guilt or shame.} \\
\texttt{* The fact that Eva is accusing Nox of hiding in his studio, which could imply that Nox is avoiding confrontation or responsibility, and this could be a source of guilt or shame for him.} \\
\texttt{* The tone of Eva's accusations, which could be perceived as condescending or mocking, and this could further trigger Nox's emotional response.} \\
\texttt{* I will also consider the fact that Nox's statement "Not after everything" suggests a past conflict or trauma, which could be a significant factor in his emotional response.} \\
\texttt{* I will analyze the emotional tone of Nox's statement "I'm just trying to move forward" and consider whether it is a genuine attempt to move on or a defensive mechanism to avoid confronting his past.} \\
\texttt{* Additionally, I will consider the fact that Eva's accusations are not just about Nox's actions, but also about his character and his past behavior, which could be a source of shame and guilt for him.} \\
\texttt{* I will also analyze the power dynamics in the conversation, and how Eva's accusations are likely to affect Nox's emotional state, given their past relationship and the context of the dialogue.} \\
\\ 
\texttt{By considering these factors, I will aim to provide a more nuanced and accurate assessment of Nox's emotional state.} \\
            }
        }
    }
    & \\ \\

\end{supertabular}
}

\end{document}
